\subsection{Tate pairs compatible with the common first skeleton}
\label{arch136:sk1-tate}

\paragraph{Archive scope and input.}
The conditions are cote136, physical pp.~117 and119, footers116
and118, scans117 and119; the common-source question is physical
p.~121, footer120, scan121. Use the actual replacements of
AQ-136-FIRST-SKELETON, whose common truncated theory is denoted
by \(F\). Thus \(F_0=S=\Gamma_{\mathbb Z}(T)\), and \(F_1\)
is the displayed free algebra on even zero-form unit disks.
All new arguments below are elementary, tagged none.

\paragraph{ARCH-136-SKELETON-TATE-LIFTS.}
Let \(f_j:A\to C_j\) be maps with \(A_0=S\), where the
admissible targets have the same first skeleton \(F\), and all
\(f_j\) agree on that skeleton. This means that the truncated
maps \(\operatorname{sk}_1 A\to F\) are identical; it does
not require \(A_1=F_1\). Fix \(q\geq2\). In a sufficiently
large weight \(w\), uniformly in the simplex and in \(N\geq1\),
suppose
\[
 z\in A_n^{w,q},\quad b\in A_n^{w,q-1},\qquad
 dz=0,\quad db=Nz.                                      \tag{ST.1}
\]
For \(q=2\) and for every \(q\geq4\), the free Tate pair
\[
 dc=z,\qquad dt=Nc-b,
 \qquad |c|=q-1,\quad |t|=q-2                           \tag{ST.2}
\]
can be mapped to all targets so that the extended maps still
agree on their first skeleton and simplex zero remains \(S\).
For \(q=3\), the exact additional obstruction for this pair is
a difference class in
\[
                  H^1(\operatorname{sk}_1\Delta[n];S^w/NS^w).
                                                               \tag{ST.3}
\]
Its vanishing for every pair of targets is necessary and
sufficient. This is a statement about the prescribed attachment,
not a vanishing theorem for (ST.3).

\begin{proof}
Put \(G=\operatorname{sk}_1\Delta[n]\). Choose the weight
bound above the common endpoint-filling bound of \(F\), the
targets' vertical comparison bounds through degree \(q\),
and their row bounds through degree \(q-1\).
The graph-incidence calculation in AQ-136-FIRST-SKELETON gives
\[
 H^0(F(G)^{w,*})=S^w,\qquad
 H^1(F(G)^{w,*})=H^1(G;S^w),\qquad H^{\geq2}=0.          \tag{ST.4}
\]
The group in degree one is free abelian. Each target on the
whole simplex has coefficient cohomology and no positive
cohomology through the required range.

For any complex \(K\) use
\[
 P_N(K)^r=K^r\oplus K^{r+1},\qquad
 d(x,y)=(dx-Ny,-dy).
\]
The usual cone exact sequence is valid without module flatness.
By (ST.4), its graph cohomology is
\[
 \begin{aligned}
 H^0(P_N(F(G)))&=S^w/NS^w,\\
 H^1(P_N(F(G)))&=H^1(G;S^w/NS^w),\qquad H^{r\geq2}=0.
 \end{aligned}
                                                               \tag{ST.5}
\]
For each whole target simplex the corresponding \(H^0\) is
\(S^w/NS^w\) and the positive groups needed here vanish.
These are canonical identifications, induced by coefficients.

Vertical exactness gives target primitives \(c_j,t_j\) as in
(ST.2). The element \((b,z)\) is a cone cycle in degree
\(q-1\), and \(y_j=(-t_j,-c_j)\) is a cone primitive of it.
All input data agree on \(G\), so the differences
\(y_j|_G-y_1|_G\) are cone cycles in degree \(q-2\) of the
same graph algebra. If \(q\geq4\), (ST.5) makes every such
difference a boundary. Choose its cone primitive on \(G\).
Both of its components extend to the target simplex: row
acyclicity fills the relative free normalized simplex chains,
one basis element at a time. Subtracting the differential of
that extension makes the restrictions of \(y_j\) and \(y_1\)
identical, without changing their differential.

For \(q=2\), first subtract a constant coefficient cone cycle
representing the degree-zero difference class in (ST.5).
That cycle extends to the whole simplex. The remaining difference
is exact. Its primitive is in cone degree minus one, whose sole
component is a zero-form; the same row-filling argument extends
it and makes the restrictions identical. Equivalently, the
difference of first primitives on the graph has a degree-one
period annihilated by \(N\); freeness in (ST.4) makes that period
zero, and the remaining second-primitive difference is a constant.

For \(q=3\), (ST.5) gives exactly (ST.3). Changing a whole
target primitive adds a degree-one cone cycle. Its cohomology
class is zero in that target, so its restriction changes the
graph difference by a boundary. Thus the stated difference
class is independent of all chosen primitives. Necessity of
its vanishing follows from equal graph restrictions. If it
vanishes, the same extension of a cone primitive used above
proves sufficiency.

Finally attach the generators parametrized by \(\mathbb Z[\Delta[n]]\),
with their vertex values identified with those just chosen.
In positive form degrees the vertex values are zero. For
\(q=2\), the zero-form \(t\) has the common coefficient
values supplied above. The differential preserves these
relations, hence also their full PD ideal. The free strict PD
universal property gives the extension and all target maps.
Its new degree-one-simplex generators have identical images in
\(F_1\), by construction, so first-skeleton agreement persists.
The vertex relations leave simplex zero unchanged. In each
simplex the vertex inclusion has a free complement, and over
\(\mathbb Q\) the change \(c'=c-b/N\) makes each new pair
a free unit disk. This also verifies preservation of the
rational vertical comparison. No integral exactness is inferred.
\end{proof}

\paragraph{Routine steps preserve the same agreement.}
The free horizontal-cell steps of AQ-136-PARTIAL-SOURCE can
also be chosen to preserve it. Cells whose new normalized
dimension is at least two add no edges. For a new edge, its
boundary is in \(A_0=S\); a positive-form boundary is zero
and can use the common zero filler, while a zero-form coefficient
in weight \(w\) has the common filler supplied by \(F_1\)
once the row cutoff is above the endpoint-filling bound.
The differential partner has the same common image.

AQ-136-RELATIVE-H1 preserves the agreement whenever its other
hypotheses hold. In normalized simplicial degree one, an old
one-cycle has the same image in the targets' common \(F_1\).
Its zero-form primitive in the normalized \(F_1\) is unique,
because both normalized vertical \(H^0\) and \(H^1\) vanish.
The extended zero-form generator therefore has the same image
in every target. Normalized degree zero is the coefficient
identity; Dold--Kan gives the desired agreement on all edges.
The strict PD pushout and its filtered iteration retain it.

Thus degree-three Tate pairs are the sole unresolved step
among these specified operations for retaining the common
first-skeleton maps. This statement supplies no automatic
vanishing of their graph residue. The already checked finite
skeleton counterexample with \(d=2\) shows that it can be
nonzero for a prescribed source and maps.
